\qhead{6. Communicating Value and Prices. How does your company communicate value? What factors are most important in your company's communication of prices?}

BYD's communication rests on the psychology of how people read a price, not just on the number itself \cite{nagle2018}.

The clearest example is the way it moves the reference point. ``Electricity cheaper than oil'' invites the buyer to compare the Qin not with other electric cars but with a petrol one, and against a Corolla the Qin looks like a bargain and an upgrade at the same time. BYD simply chooses the comparison that makes the surplus look largest. The fuel-saving story works the same trick from a different angle. Framing the roughly \euro1{,}240 a year as money the owner stops losing, rather than as a feature to be gained, leans on the fact that losses feel heavier than equivalent gains, so the message lands harder.

Giving the ``God's Eye'' driver assistance away for free, while Tesla charges thousands for its self-driving package \cite{insidechinaADAS}, does two things. It knocks out the rival's price anchor, and it avoids the problem that bolting a costly option onto a cheap car feels proportionally enormous, whereas folding it into the price hides that shock.

The factor that matters most, though, is the inference people draw between price and quality. In Europe a car priced at under half a Tesla risks being read as cheap in the bad sense, and getting that wrong would turn a price advantage into a brand problem. A car is really an experience good that carries both economic and psychological value, which is the value-communication grid the course sets out, so BYD's task splits in two. On the economic side it offers assurance the buyer can weigh in advance, the Blade Battery safety results, the warranties, the awards and test drives. On the psychological side it offers assurance through brand signals, the Yangwang halo and a Hungarian factory that shows it is in Europe to stay. The thread running through all of it is the same message: the low price reflects a real cost advantage, not a compromise on quality. The ladder helps carry that message, because putting the Yangwang halo cars in view first anchors expectations high, so a mid-range BYD then reads as sensible value rather than a cheap option.
